\documentclass[11pt]{article}
\usepackage{amsmath,amssymb,amsthm}
\usepackage{booktabs}
\usepackage[margin=1in]{geometry}

\newtheorem{proposition}{Proposition}

\title{Jukes--Cantor Distance for the 6502life Noisy Copy Channel}
\author{6502life project}
\date{March 2026}

\begin{document}
\maketitle

\begin{abstract}
We derive a closed-form evolutionary distance formula for cells diverging
through the 6502life noisy copy channel, where each bit is independently
randomized with probability $\varepsilon$ per copy event. This is the
binary Jukes--Cantor model. The MLE of evolutionary distance is a
closed-form function of Hamming distance.
\end{abstract}

\section{The Noisy Copy Channel}

When a BRK noisy-copy fires (operands \$F5--\$F8), the origin cell's
$M = 1024$ bytes are copied to a neighbor. Each bit is independently
randomized (replaced by a fair coin flip) with probability
$\varepsilon$ (\texttt{pBitNoise}), and faithfully copied with
probability $1-\varepsilon$. The default is $\varepsilon = 1/2048$,
giving $\sim$1 bit error per 256-byte page copied.

\section{Binary Jukes--Cantor Distance}

\begin{proposition}
After $T$ copy events on the tree path between two cells, the fraction
of differing bits is
\begin{equation}
p = \tfrac{1}{2}\bigl(1 - (1-\varepsilon)^T\bigr)
\end{equation}
and the MLE of $T$ given observed bit mismatch fraction $\hat{p}$ is
\begin{equation}
\boxed{\hat{T} = \frac{\log(1 - 2\hat{p})}{\log(1 - \varepsilon)}}
\end{equation}
Saturation: $p \to 1/2$ as $T \to \infty$.
\end{proposition}

\begin{proof}
Each bit evolves under a 2-state symmetric Markov chain with transition
rate $\varepsilon/2$ per copy event. Starting from ``same'' at $T=0$:
$P(\text{diff}) = \frac{1}{2}(1 - (1-\varepsilon)^T)$. This is the
standard binary Jukes--Cantor model. The MLE of $p$ from $n$ i.i.d.\
bit comparisons is $\hat{p} = (\text{differing bits})/n$, and the
invariance property of MLEs gives the formula for $\hat{T}$.
\end{proof}

\subsection{Byte-level distance}

Since bits within a byte are independent (no byte-level noise), the
probability that two corresponding bytes match is $(1-p)^8$:
\begin{align}
P(\text{byte differs}) &= 1 - \Bigl[\frac{1 + (1-\varepsilon)^T}{2}\Bigr]^8
\end{align}
Inverting:
\begin{equation}
\hat{T} = \frac{\log\bigl(2(1-\hat{p}_{\text{byte}})^{1/8} - 1\bigr)}{\log(1-\varepsilon)}
\end{equation}

\subsection{Byte Hamming histogram}

The per-byte Hamming distance $h \in \{0,\ldots,8\}$ follows
$\mathrm{Bin}(8, p)$. The sufficient statistic is the sample mean
$\bar{h}$, giving $\hat{p} = \bar{h}/8$. The histogram provides no
additional information beyond the mean (it is a single-parameter
exponential family), but can be used for goodness-of-fit testing
against the copy-noise model.

\section{Applicability}

The distance formula assumes all sequence divergence arises from the
noisy copy channel. In practice, cross-contamination (random cells
writing into each other's code regions) is often the dominant source
of divergence and is not phylogenetically structured. The byte Hamming
histogram can diagnose this: under the copy-noise model, $h \sim
\mathrm{Bin}(8,p)$; excess weight at $h \approx 4$ relative to
$h = 0$ indicates non-copy contamination.

\end{document}
